% =====================================================================
%  introduction.tex
%  Introduction chapter for a thesis on imputation of missing data in
%  sparse financial panels, evaluated by downstream model quality.
%
%  USAGE: This is a chapter FRAGMENT. Do NOT add \documentclass or a
%  preamble here. Include it from your report-class main.tex, e.g.
%
%      \documentclass[12pt]{report}
%      \usepackage{amsmath,amssymb}     % needed for the math below
%      % ... your bibliography setup (natbib / biblatex / plain bibtex)
%      \begin{document}
%      \include{introduction}           % or \input{introduction.tex}
%      ...
%      \end{document}
%
%  Citations use the package-agnostic \cite{} so this compiles under
%  natbib, biblatex, or plain BibTeX. The keys used here are listed at
%  the end of this file; matching .bib entries are provided separately.
% =====================================================================

\chapter{Introduction}
\label{chap:introduction}

\section{Motivation}
\label{sec:motivation}

Financial time series data is never fully observed.
Limit-order book feeds drop quotes during bursts of order cancellations.
Equity markets suspend trading when prices move too violently.
Over-the-counter bond markets leave entire yield-curve tenors unquoted
for days.
Implied volatility surfaces contain large regions where no liquid
option trade has occurred.
Macro and alternative data arrive at mixed, irregular frequencies, with
publication lags that create structural gaps at the daily level.
This missing data is pervasive and often elusive, and gets in the way
of advanced analysis. Thus practitioners or a researchers face
the same decision: what value should be attributed to the missing
observation, and how much uncertainty should be attached to that
attribution?

The simplest and most common solution is to carry forward the last observed value,
but this can lead to systematic mispricing, underestimation of risk, and distortion
of correlations among assets.

We aim to take a different path, one that looks for solutions that aim to preserve as much
as possible the statistical properties of the series, in order to maintain the integrity of
the underlying data and consequently the robustness of any downstream model that relies on it.
Since this can be considered an inference problem, we will also explore ML algorithms,
and in classical machine learning approach, we will
focus on quality of inference rather than interpretability of the model.

In this thesis we will test these different methods of imputation, and evaluate them not only on the
accuracy of the imputation itself (MSE,...),
but also on the quality of the downstream model that uses the imputed data, and how
far it deviates from the case where all data is observed.

\section{Literature review}
\label{sec:literature}

The modern theory of missing data originates with
\cite{rubin1976}, who classifies the mechanism that generates the gaps.
Let the complete data be partitioned into observed and missing parts,
$X = (X_{\text{obs}}, X_{\text{mis}})$, and let $R$ denote the matrix of
missingness indicators, governed by a conditional law $P(R \mid X,
\phi)$. The mechanism is \emph{missing completely at random} (MCAR) if
$R$ is independent of the data, \emph{missing at random} (MAR) if it
depends only on the observed part $X_{\text{obs}}$, and \emph{missing not
at random} (MNAR) if it depends on the unobserved values
$X_{\text{mis}}$ themselves. The practical importance of this taxonomy
lies in the \emph{ignorability} result: likelihood-based and Bayesian
inference may ignore the missingness mechanism if and only if the data
are MAR and the data and mechanism parameters are distinct
\cite{rubin1976,littlerubin2019}. This is the condition that licenses
the expectation--maximisation algorithm \cite{dempster1977} and multiple
imputation \cite{rubin1987} without explicitly modelling $R$.

Financial data violates the convenient end of this taxonomy. Missingness
in returns and fundamentals is typically MNAR and strongly structured:
it clusters in the cross section and persists over time, and the very
fact that a value is missing is informative about that value. Returns
themselves depend on whether data are missing
\cite{bryzgalova2025}. Two consequences follow directly and motivate the
methods studied here. First, complete-case analysis conditions on
survival and on disclosure, producing a sample that is not
representative of the population of interest. Second, unconditional mean
imputation, even where it leaves first moments unbiased, attenuates
estimated variances and covariances toward zero. Because mean--variance
portfolio selection and risk measurement consume the covariance matrix
$\Sigma$---and, through the optimiser, its inverse $\Sigma^{-1}$---this
attenuation corrupts exactly the object the downstream model depends
upon.

\section{From imputation to estimation: the downstream view}
\label{sec:downstream}

The clearest illustration of the downstream view, and the running
example of this thesis, is mean--variance portfolio selection
\cite{markowitz1952}, in which the optimal risky portfolio is
proportional to $\Sigma^{-1}\mu$ for a mean vector $\mu$ and covariance
$\Sigma$. It is well established that plugging sample estimates into this
expression performs poorly out of sample: the optimiser behaves as an
``error maximiser'', loading on assets whose moments are most favourably
mis-estimated \cite{michaud1989,jobsonkorkie1981}. The estimation-error
problem is severe enough that a naive equally weighted portfolio is hard
to beat \cite{demiguel2009}, and a large literature responds with
regularisation: shrinkage of the covariance matrix
\cite{ledoitwolf2004,ledoitwolf2017} or equivalent weight constraints
\cite{jagannathanma2003}.

Imputation enters this picture not as a separate preprocessing step but
as a further source of regularisation. Filling a return panel with a
low-rank or factor-based reconstruction shrinks the estimated covariance
toward a factor structure---the same target used deliberately by
shrinkage estimators. The choice of imputation method and the choice of
covariance estimator are therefore not separable, and the error that
imputation introduces propagates through $\Sigma^{-1}$ into the portfolio
weights and ultimately into realised performance. This is why
reconstruction accuracy is an inadequate criterion: it ignores how the
imputation interacts with the conditioning of the downstream estimator.
The natural evaluation is instead a two-layer one, comparing methods both
on reconstruction error and on a downstream objective such as realised
portfolio variance, Sharpe ratio, or tail risk, with particular interest
in the configurations where the rankings reverse.

Mean--variance optimisation is only the leading example. The same
impute-then-estimate logic governs risk-model estimation and the
computation of Value-at-Risk and Expected Shortfall; the estimation of
factor models and cross-sectional return predictors from panels of firm
characteristics \cite{freyberger2020,chenzimmermann2022}; volatility
forecasting; macroeconomic nowcasting from mixed-frequency, ragged-edge
data; the construction of yield curves and implied-volatility surfaces
from sparsely observed grids; and credit-risk and ESG scoring, where the
underlying data are both heavily and informatively missing. In each case
the relevant question is identical: how does the imputation choice change
the estimate that the practitioner actually uses.

\section{Sparse and illiquid markets}
\label{sec:emerging}

The problem sharpens in emerging and illiquid markets, where it also
changes in character. The dominant difficulty there is often not an
explicit missing cell but \emph{staleness}: an asset that does not trade
still carries a recorded closing price---yesterday's---so the panel can
appear fully observed while being silently contaminated by repeated
prices and spurious zero returns. The missingness is in part
\emph{latent}, inferred rather than flagged, which must be detected
before it can be addressed, for instance through the proportion of
zero-return days and related liquidity proxies
\cite{lesmond1999,amihud2002}. Whether a stale observation should be
treated as observed or converted to missing is itself a modelling
decision with downstream consequences.

Illiquidity also intensifies the statistical biases of interest.
Non-trading is more likely in adverse states of the world, so the
missing observations are disproportionately the extreme ones, which
inflates estimated means and understates downside risk. Non-synchronous
prices bias correlations and betas downward and understate volatility, so
that a portfolio optimiser fed stale data systematically understates risk
and overstates diversification; standard corrections estimate exposures
using leads and lags \cite{scholeswilliams1977,dimson1979}. Because the
characteristic failure mode is the under-statement of risk, evaluation in
this setting should emphasise realised downside measures rather than
in-sample Sharpe ratios alone. Finally, illiquidity yields a testable
prediction that the thesis examines: since the time series of an illiquid
asset is largely uninformative, cross-sectional and low-rank methods,
which borrow strength from contemporaneous liquid peers, should outperform
time-series fillers by a wider margin in emerging markets than in
developed ones.

\section{Research questions}
\label{sec:questions}


\section{Contributions}
\label{sec:contributions}


\section{Outline of the thesis}
\label{sec:outline}